\clearpage\pagestyle{otMain}\setlength{\otBlocEspacement}{2.5mm}\fontsize{11}{15}\selectfont\corrigesChapitre{20}{Lire et croiser des données}
\begin{blocLivre}{cor-20-1}\begin{corrige}{1}{Le tableau de l’enquête}\demarche{Contrôler le total avant de conclure.}
Total : $12+9+5+4=\res{30}$ élèves. Le $\res{\text{foot}}$ arrive en tête avec 12 voix.
\end{corrige}\end{blocLivre}
\begin{blocLivre}{cor-20-2}\begin{corrige}{2}{Double entrée}\demarche{Croiser la ligne et la colonne demandées.}
Croisement CM2 B-forêt : $\res{17}$. Colonne musée : $15+9=\res{24}$ élèves.
\end{corrige}\end{blocLivre}
\begin{blocLivre}{cor-20-3}\begin{corrige}{3}{Lire les barres}\demarche{Multiplier les graduations par leur pas.}
Lundi : $\res{15}$ ; jeudi : $\res{35}$, la barre la plus haute : record le $\res{\text{jeudi}}$.
\label{fig-cor-20-3}\figureGeo[corrige]{don-3}
\end{corrige}\end{blocLivre}
\begin{blocLivre}{cor-20-4}\begin{corrige}{4}{La courbe de l’année}\demarche{Chercher les points extrêmes.}
Le point le plus bas est en $\res{\text{janvier}}$ ($3\,^\circ$C) ; juillet culmine à $23\,^\circ$C : écart $23-3=\res{20}$ degrés.
\label{fig-cor-20-4}\figureGeo[corrige]{don-4}
\end{corrige}\end{blocLivre}
\begin{blocLivre}{cor-20-5}\begin{corrige}{5}{Le circulaire de la cantine}\demarche{Traduire chaque secteur en fraction du total.}
A : $60:2=\res{30}$ ; B : $60:4=\res{15}$ ; C : le reste, $60-30-15=\res{15}$ élèves.
\label{fig-cor-20-5}\figureGeo[corrige]{don-5}
\end{corrige}\end{blocLivre}
\begin{blocLivre}{cor-20-6}\begin{corrige}{6}{Construire le diagramme}\demarche{Une barre par sport, hauteur égale à l’effectif.}
Quatre barres de hauteurs 12, 9, 5 et 4 : le dessin figure en vert, chaque hauteur lue sur l’axe gradué de 2 en 2.
\label{fig-cor-20-6}\figureGeo[corrige]{don-6}
\end{corrige}\end{blocLivre}
\begin{blocLivre}{cor-20-7}\begin{corrige}{7}{Compléter le tableau}\demarche{Utiliser chaque total de ligne.}
Mardi-entrées : $313-145=\res{168}$. Lundi-goûters : $184-96=\res{88}$. Les totaux pilotent les cases manquantes.
\end{corrige}\end{blocLivre}
\begin{blocLivre}{cor-20-8}\begin{corrige}{8}{Les points du repère}\demarche{Comparer les écarts entre points voisins.}
Semaine 2 : $\res{14}$ cm ; semaine 5 : $\res{38}$ cm. Le plus grand bond sépare les semaines $\res{3\text{ et }4}$ : de 18 à 30 cm, soit 12 cm.
\label{fig-cor-20-8}\figureGeo[corrige]{don-8}
\end{corrige}\end{blocLivre}
\begin{blocLivre}{cor-20-9}\begin{corrige}{9}{Vrai ou faux sur le circulaire}\demarche{Chercher le total avant de convertir.}
$\res{\text{Non}}$ sans connaître le total : un quart de 60 fait 15, un quart de 100 fait 25. Un secteur donne une part du tout, pas un effectif tout seul.
\end{corrige}\end{blocLivre}
\begin{blocLivre}{cor-20-10}\begin{corrige}{10}{Deux documents à croiser}\demarche{Prendre une donnée dans chaque document.}
$76-48=\res{28}$ pains l’après-midi. Croiser deux documents, c’est soustraire ou additionner leurs informations.
\label{fig-cor-20-10}\figureGeo[corrige]{don-10}
\end{corrige}\end{blocLivre}
\begin{blocLivre}{cor-20-11}\begin{corrige}{11}{La barre incohérente}\demarche{Comparer chaque barre à sa valeur.}
La barre $\res{\text{CE1}}$ monte à 22 au lieu de $\res{26}$ ; la correction figure en vert. Toujours confronter graphique et données sources.
\label{fig-cor-20-11}\figureGeo[corrige]{don-11}
\end{corrige}\end{blocLivre}
\begin{blocLivre}{cor-20-12}\begin{corrige}{12}{L’enquête à construire}\demarche{Compter en barrant les réponses une à une.}
Effectifs : 0 livre : $\res3$ ; 1 livre : $\res5$ ; 2 livres : $\res3$ ; 3 livres : $\res1$. Total 12 réponses. Au moins 2 livres : $3+1=\res4$ élèves.
\end{corrige}\end{blocLivre}
\begin{blocLivre}{cor-20-13}\begin{corrige}{13}{Le climat en deux graphiques}\demarche{Lire chaque document puis croiser.}
La courbe culmine en $\res{\text{juillet}}$ ($23\,^\circ$C) et la barre de juillet est la plus basse (18 mm) : juillet cumule les deux records. Croiser les deux documents répond à la question.
\label{fig-cor-20-13}\figureGeo[corrige]{don-13}
\end{corrige}\end{blocLivre}
\begin{blocLivre}{cor-20-14}\begin{corrige}{14}{Choisir le bon graphique}\demarche{Associer question et point fort de chaque outil.}
L’évolution : la $\res{\text{courbe}}$. Les parts du tout : le $\res{\text{circulaire}}$. La valeur exacte : le $\res{\text{tableau}}$. Chaque question a son graphique.
\end{corrige}\end{blocLivre}
\begin{blocLivre}{cor-20-15}\begin{corrige}{15}{Le piège de la lecture rapide}\demarche{Vérifier l’origine de l’axe avant de comparer.}
L’axe tronqué ne montre que ce qui dépasse 50 : les hauteurs dessinées (10 contre 20) doublent, mais les ventes réelles passent seulement de 60 à $\res{70}$. Lire les valeurs, pas l’impression visuelle.
\end{corrige}\end{blocLivre}
